%% kapa: exam=1 round=23 year=2012 period=2 form=A subject=accounting subject_ko=회계학 seq=31 number=31
%% source: Q-Net 감정평가사 기출문제 (cst00309) artlSeq=5052580 file=raw/1008/5052580/2012년 제23회 감정평가사 제1차시험 기출문제(2교시 A형).pdf
%% conversion: kordoc 4.24.51 -> tools/gen1.py (auto); review=proofread-sonnet; flags=none
\begin{question}{31}
제조부문 A, B와 보조부문 X, Y의 서비스 제공관계는 다음과 같다.
\begin{center}
\small
\begin{tabular}{|c|c|c|c|c|c|}
\hline
 & \multicolumn{2}{c|}{보조부문} & \multicolumn{2}{c|}{제조부문} & \multirow{2}{*}{합 계} \\
\cline{2-5}
 & X & Y & A & B & \\
\hline
X & － & 40단위 & 20단위 & 40단위 & 100단위 \\
\hline
Y & 80단위 & － & 60단위 & 60단위 & 200단위 \\
\hline
\end{tabular}
\end{center}
X, Y부문의 원가는 각각 ￦160,000, ￦200,000이다. 단계배부법에 의해 X부문을 먼저 배부하는 경우와 Y부문을 먼저 배부하는 경우의 제조부문 A에 배부되는 총보조부문원가의 차이는? \par
\begin{choices}
  \item ￦24,000
  \item ￦25,000
  \item ￦26,000
  \item ￦27,000
  \item ￦28,000
\end{choices}
\end{question}
\answer{1}
